\input{preamble.tex}

\title{Лабораторная работа № 2. \\ Предварительная настройка оборудования Cisco}
\author{Сущенко Алина \\ НПИбд-01-23}

\institute{Российский университет дружбы народов имени Патриса Лумумбы}
\date{2026}

\begin{document}

\frame{\titlepage}

\begin{frame}
\frametitle{Цель работы}
\begin{itemize}
    \item Получить основные навыки по начальному конфигурированию оборудования Cisco.
\end{itemize}
\end{frame}

\begin{frame}[containsverbatim]
\frametitle{Выполнение работы}
В логической рабочей области \texttt{Packet Tracer} разместим коммутатор, маршрутизатор и 2 оконечных устройства типа PC (PC0 и PC1), соединим один PC с маршрутизатором, другой PC — с коммутатором, как показано на схеме.
    \centering
    \includegraphics[width=\textwidth]{../images/01.png}
    \captionof{figure}{Схема подключения оборудования для проведения его предварительной настройки.}
\end{frame}

\begin{frame}
\frametitle{Конфигурация маршрутизатора}
    \centering
    \includegraphics[width=\textwidth]{../images/02.png}
    \captionof{figure}{Настройка маршрутизатора.}
\end{frame}

\begin{frame}
\frametitle{Конфигурация коммутатора}
    \centering
    \includegraphics[width=\textwidth]{../images/03.png}
    \captionof{figure}{Настройка коммутатора.}
\end{frame}

\begin{frame}
\frametitle{Проверка PC0}
    \centering
    \includegraphics[width=\textwidth]{../images/04.png}
    \captionof{figure}{Просмотр текущих настроек PC0.}
\end{frame}

\begin{frame}
\frametitle{Проверка PC1}
    \centering
    \includegraphics[width=\textwidth]{../images/05.png}
    \captionof{figure}{Просмотр текущих настроек PC1.}
\end{frame}

\begin{frame}[containsverbatim]
\frametitle{Проверка работоспособности соединений}
Теперь проверим связь при помощи наших PC. На PC0 выполним команды:
\begin{minted}{bash}
ping 192.168.2.10
telnet 192.168.2.1
ssh -l admin 192.168.2.1
\end{minted}
    \centering
    \includegraphics[width=\textwidth]{../images/06.png}
    \captionof{figure}{Проверка работоспособности для PC0.}
\end{frame}

\begin{frame}[containsverbatim]
\frametitle{Проверка работоспособности соединений}
Затем то же самое, но для PC1:
\begin{minted}{bash}
ping 192.168.1.254
telnet 192.168.1.254
ssh -l admin 192.168.1.254
\end{minted}
    \centering
    \includegraphics[width=\textwidth]{../images/07.png}
    \captionof{figure}{Проверка работоспособности для PC1.}
\end{frame}

\begin{frame}
\frametitle{Контрольные вопросы}
\begin{block}{1. Какие существуют способы подключения к сетевому оборудованию?}
Существуют следующие способы подключения к сетевому оборудованию:
\begin{itemize}
    \item \textbf{Консольное подключение (Console)} — через порт RS-232 с использованием консольного кабеля (синий кабель в Packet Tracer). Требуется для первоначальной настройки, не требует сетевых настроек.
    \item \textbf{Telnet} — удалённое подключение по протоколу Telnet (порт 23). Работает поверх IP-сети, но данные передаются в незашифрованном виде.
    \item \textbf{SSH} — защищённое удалённое подключение (порт 22). Использует шифрование, требует настройки домена и генерации RSA-ключей.
    \item \textbf{AUX-порт} — дополнительный порт для подключения через модем (используется редко).
\end{itemize}
\end{block}
\end{frame}

\begin{frame}
\frametitle{Контрольные вопросы}
\begin{block}{2. Какие режимы работы командной строки существуют в оборудовании Cisco?}
В оборудовании Cisco существуют следующие основные режимы:
\begin{itemize}
    \item \textbf{Пользовательский режим (User EXEC Mode)} — режим с ограниченными правами, обозначается символом \texttt{>} (например, \texttt{Router>}). Доступны только команды для просмотра базовой информации.
    \item \textbf{Привилегированный режим (Privileged EXEC Mode)} — режим с полным доступом к командам просмотра и отладки, обозначается символом \texttt{#} (например, \texttt{Router#}). Вход осуществляется командой \texttt{enable}.
    \item \textbf{Режим глобальной конфигурации (Global Configuration Mode)} — режим для изменения конфигурации устройства, обозначается \texttt{(config)\#}. Вход осуществляется командой \texttt{configure terminal}.
    \item \textbf{Режимы конфигурации интерфейсов} — подрежимы для настройки конкретных интерфейсов (\texttt{(config-if)\#}).
    \item \textbf{Режимы конфигурации линий} — подрежимы для настройки консольного порта и VTY-линий (\texttt{(config-line)\#}).
\end{itemize}
\end{block}
\end{frame}

\begin{frame}
\frametitle{Контрольные вопросы}
\begin{block}{3. Какие команды используются для просмотра конфигурации оборудования?}
Для просмотра конфигурации используются следующие команды:
\begin{itemize}
    \item \texttt{show running-config} — просмотр текущей конфигурации в оперативной памяти (изменённой, но не сохранённой).
    \item \texttt{show startup-config} — просмотр конфигурации, сохранённой в NVRAM (загружается при перезагрузке).
    \item \texttt{show ip interface brief} — краткая информация по всем интерфейсам: IP-адреса, статус, протокол.
    \item \texttt{show interfaces} — детальная информация по всем интерфейсам.
    \item \texttt{show interfaces fastEthernet 0/1} — детальная информация по конкретному интерфейсу.
    \item \texttt{show vlan brief} — просмотр информации о VLAN на коммутаторе.
    \item \texttt{show ip route} — просмотр таблицы маршрутизации.
    \item \texttt{show version} — информация о версии IOS и оборудовании.
\end{itemize}
\end{block}
\end{frame}

\begin{frame}
\frametitle{Контрольные вопросы}
\begin{block}{4. Каково назначение команд \texttt{no shutdown} и \texttt{shutdown}?}
\begin{itemize}
    \item \texttt{shutdown} — административно отключает интерфейс. Порт перестаёт передавать данные, статус становится \texttt{administratively down}.
    \item \texttt{no shutdown} — включает интерфейс. Порт активируется и начинает передавать данные (при наличии физического подключения статус становится \texttt{up}).
\end{itemize}
Эти команды выполняются в режиме конфигурации интерфейса (\texttt{interface}).
\end{block}
\end{frame}

\begin{frame}
\frametitle{Контрольные вопросы}
\begin{block}{5. Что такое VLAN и для чего он используется?}
\textbf{VLAN (Virtual Local Area Network)} — технология виртуальных локальных сетей, позволяющая логически разделять устройства в одной физической сети на изолированные группы.

Основные назначения VLAN:
\begin{itemize}
    \item \textbf{Сегментация сети} — разделение на логические сегменты (например, отделы компании).
    \item \textbf{Безопасность} — изоляция трафика между разными VLAN.
    \item \textbf{Уменьшение широковещательного домена} — широковещательный трафик не выходит за пределы своего VLAN.
    \item \textbf{Гибкость} — устройства из одного VLAN могут находиться на разных коммутаторах.
\end{itemize}
В лабораторной работе был создан \textbf{VLAN 2} для подсети \texttt{192.168.2.0/24}, в которую вошли порт Fa0/1 (к маршрутизатору) и порт Fa0/2 (к PC1).
\end{block}
\end{frame}

\begin{frame}
\frametitle{Контрольные вопросы}
\begin{block}{6. Как настроить доступ к оборудованию по протоколу SSH?}
Для настройки SSH необходимо выполнить следующие шаги:
\begin{enumerate}
    \item Задать имя устройства: \texttt{hostname имя}
    \item Задать доменное имя: \texttt{ip domain-name имя.домена}
    \item Сгенерировать RSA-ключи: \texttt{crypto key generate rsa}
    \item Создать локального пользователя: \texttt{username admin privilege 1 secret пароль}
    \item Настроить линии VTY: \texttt{line vty 0 4}
    \item Включить только SSH: \texttt{transport input ssh}
    \item Указать метод аутентификации: \texttt{login local}
\end{enumerate}
\end{block}
\end{frame}

\begin{frame}
\frametitle{Контрольные вопросы}
\begin{block}{7. В чём различие между Telnet и SSH?}
\begin{table}[h]
\centering
\begin{tabular}{|p{3cm}|p{5cm}|}
\hline
\textbf{Telnet} & \textbf{SSH} \\
\hline
Порт 23 & Порт 22 \\
\hline
Не шифрует данные (передаёт в открытом виде, включая пароли) & Шифрует все передаваемые данные \\
\hline
Уязвим к перехвату трафика & Защищён от перехвата \\
\hline
Не требует генерации ключей & Требует генерации RSA-ключей \\
\hline
Проще в настройке & Сложнее в настройке, но безопаснее \\
\hline
\end{tabular}
\end{table}
\end{block}
\end{frame}

\begin{frame}
\frametitle{Контрольные вопросы}
\begin{block}{8. Для чего используется команда \texttt{ping}?}
Команда \texttt{ping} (Packet Internet Groper) используется для проверки сетевой связности между устройствами. Она отправляет ICMP-запросы (Echo Request) и ожидает ICMP-ответы (Echo Reply).

Позволяет:
\begin{itemize}
    \item Проверить доступность удалённого узла.
    \item Оценить время прохождения пакета (RTT — Round Trip Time).
    \item Выявить потерю пакетов.
    \item Диагностировать проблемы с маршрутизацией и подключением.
\end{itemize}
Пример вывода успешного пинга: \texttt{Reply from 192.168.1.254: bytes=32 time<1ms TTL=255}
\end{block}
\end{frame}

\begin{frame}
\frametitle{Вывод}
В результате выполнения лабораторной работы были настроены маршрутизатор и коммутатор Cisco в среде Packet Tracer: 
\begin{itemize}
    \item настроены интерфейсы и VLAN 2;
    \item обеспечена связность между двумя ПК из разных подсетей;
    \item настроены пароли, Telnet и SSH для удалённого доступа.
\end{itemize}
Проверка командой \texttt{ping} подтвердила работоспособность сети, удалённое подключение выполняется успешно.
\end{frame}

\end{document}
